\begin{theorem}[单参对数读出与对角多通道读出的分辨率分岔：常数精度的必要代价]\label{thm:spg-single-parameter-log-readout-vs-diagonal-multichannel-bifurcation}
取互异素数 \(p_1,\dots,p_k\) 并令 \(b\in\{0,1\}^k\) 表示 \(k\)-比特 Gödel 寄存器。
定义单参对数读出
\[
u(b):=\sum_{i=1}^k b_i\log p_i.
\]
则存在 \(b\neq b'\) 使得
\[
\abs{u(b)-u(b')}
\ \le\
\frac{1}{2^k-1}\sum_{i=1}^k \log p_i.
\]
因此，任何仅观测单标量 \(u(b)\) 且具有误差尺度不小于上述最小间距量级的机制，都不可能对全部 \(b\in\{0,1\}^k\) 保证唯一译码；若要求无碰撞译码，则精度预算必随 \(k\) 指数增长。
\par
相对地，在推论 \ref{cor:spg-godel-axis-boundary-invertibility} 与定理 \ref{thm:spg-godel-boundary-decoding-relative-error-threshold} 的设定下，若可获得对角多通道比值
\[
r_i:=\frac{dB_{ii}}{B_0}=p_i^{2b_i},
\]
并满足相对误差模型
\(
\widetilde B_0=B_0(1+\delta_0),\ 
\widetilde B_{ii}=B_{ii}(1+\delta_i)
\)
且 \(|\delta_0|,|\delta_i|\le\varepsilon\)，则当
\[
\varepsilon<\frac{p_{\min}-1}{p_{\min}+1},
\qquad
p_{\min}:=\min_i p_i,
\]
最近整数译码
\[
\widetilde b_i:=\left\lfloor \frac{1}{2\log p_i}\log\!\left(\frac{d\,\widetilde B_{ii}}{\widetilde B_0}\right)\right\rceil
\]
可保证对所有 \(i\) 精确恢复 \(b_i\)。
从而，在“单参坐标”与“多通道对角坐标”之间存在严格分岔：前者的无碰撞译码要求指数级精度，后者可在与 \(k\) 无关的常数相对误差阈值下实现逐坐标精确译码。
\end{theorem}
\begin{proof}
把集合 \(\{u(b):b\in\{0,1\}^k\}\) 视为实轴上的 \(2^k\) 个点，且
\(
0\le u(b)\le \sum_{i=1}^k\log p_i
\)。
由鸽巢原理，区间 \([0,\sum_i\log p_i]\) 被这 \(2^k\) 个点划分出的相邻间距最小值不超过区间长度除以 \(2^k-1\)，从而存在 \(b\neq b'\) 满足所示不等式。
对单参读出的不可唯一性结论由“误差大于最小间距则发生碰撞”直接推出。
\par
多通道情形由定理 \ref{thm:spg-godel-boundary-decoding-relative-error-threshold} 直接给出。证毕。
\end{proof}

\endinput

